\chapterpage{Arrays and two pointers}{Eight problems. A read index and a write
index, or two ends closing in --- this one idea turns a great many quadratic
solutions linear.}

% ------------------------------------------------------------------
\problem{Move zeros to the end}{Easy}

Move every zero to the end, keeping the order of the non-zero elements, in
place.

\idea{A write pointer. Copy every non-zero forward to the write position, then
fill the tail with zeros.}

\begin{dia}
  \cellrowst{a}{0}{dead/0,soft/1,dead/0,soft/3,soft/12}
  \idxrow{-0.62}{4}
  \node[ptr] at (0.8,0.78) {write};
  \node[ptr] at (2.4,0.78) {read};
  \draw[flow] (0.8,0.62) -- (0.8,0.44);
  \draw[flow] (2.4,0.62) -- (2.4,0.44);
  \draw[flow] (2.4,1.15) to[out=160,in=20] (0.9,1.15);
  \caplabel{5.0}{0.78}{read runs ahead; write only advances on a non-zero};
  \cellrowst{b}{-1.7}{hot/1,hot/3,hot/12,dead/0,dead/0}
  \node[note, anchor=east] at (-0.5,-1.7) {after};
\end{dia}

\textbf{How it works.} Two indices move at different speeds over the same
array. \texttt{read} visits every slot; \texttt{write} only advances when
something is actually kept. Because \texttt{write} can never get ahead of
\texttt{read}, copying \texttt{a[read]} into \texttt{a[write]} can only ever
overwrite a slot that has already been read --- which is why this is safe to
do in place with no temporary array.

When the first loop ends, \texttt{write} is the count of non-zero elements, so
everything from there to the end is tail that must be zeroed.

\solheading
\begin{lstlisting}[style=java]
public static void moveZeros(int[] a) {
    if (a == null) return;
    int write = 0;
    for (int read = 0; read < a.length; read++)
        if (a[read] != 0) a[write++] = a[read];
    while (write < a.length) a[write++] = 0;
}
\end{lstlisting}

\cost{$O(n)$ time, $O(1)$ space.}
\trap{Swapping with the \emph{end} of the array instead of copying to
\texttt{write} destroys the relative order of the non-zero elements, which the
statement asks you to preserve.}

% ------------------------------------------------------------------
\problem{Remove duplicates from a sorted array}{Easy}

Compact a sorted array in place so each value appears once; return the new
length.

\idea{The same write pointer. Because the array is sorted, duplicates are
always adjacent, so compare each element with the last one you kept.}

\begin{dia}
  \cellrowst{a}{0}{soft/1,dead/1,soft/2,dead/2,dead/2,soft/3}
  \idxrow{-0.62}{5}
  \node[ptr] at (1.6,0.78) {write};
  \node[ptr] at (3.2,0.78) {read};
  \draw[flow] (1.6,0.62) -- (1.6,0.44);
  \draw[flow] (3.2,0.62) -- (3.2,0.44);
  \cellrowst{b}{-1.7}{hot/1,hot/2,hot/3,dead/2,dead/2,dead/3}
  \draw[span] (-0.4,-2.3) -- (1.2,-2.3);
  \caplabel{0.4}{-2.85}{returned length 3; the rest is ignored};
\end{dia}

\textbf{How it works. }Sortedness is doing the heavy lifting: it guarantees
that if a value is going to repeat, it repeats immediately. So the test for
``have I seen this already'' collapses from a set lookup to a single
comparison with the previous kept element.

That comparison must be against \texttt{a[write - 1]}, not
\texttt{a[read - 1]}. Once you start overwriting, \texttt{a[read - 1]} may be
a value that was copied there this pass, so it no longer tells you what the
last distinct value was.

\solheading
\begin{lstlisting}[style=java]
public static int dedupeSorted(int[] a) {
    if (a == null || a.length == 0) return 0;
    int write = 1;
    for (int read = 1; read < a.length; read++)
        if (a[read] != a[write - 1]) a[write++] = a[read];
    return write;
}
\end{lstlisting}

\cost{$O(n)$ time, $O(1)$ space.}
\trap{Comparing against \texttt{a[read - 1]}. It looks equivalent and passes
on inputs with no more than two consecutive duplicates.}

% ------------------------------------------------------------------
\problem{Two-sum in a sorted array}{Easy}

Return the indices of two values that add to a target, or \texttt{null}.

\idea{One pointer at each end. If the sum is too big the right pointer must come
in; if it is too small the left must go out. No other move can help.}

\begin{dia}
  \cellrowst{a}{0}{soft/2,cell/7,cell/11,cell/15,soft/19}
  \idxrow{-0.62}{4}
  \node[ptr] at (0,0.78) {i};
  \node[ptr] at (3.2,0.78) {j};
  \draw[flow] (0,0.62) -- (0,0.44);
  \draw[flow] (3.2,0.62) -- (3.2,0.44);
  \node[note, anchor=west] at (3.9,0) {$2+19=21 > 18$};
  \node[note, anchor=west] at (3.9,-1.0) {so \texttt{j--}: nothing paired};
  \node[note, anchor=west] at (3.9,-1.5) {with 19 can be small enough};
\end{dia}

\textbf{How it works.} The argument for why nothing is missed matters more
than the code. Suppose the sum at the current pair is too large. Every pair
using the right-hand value and \emph{any} remaining left-hand value is at
least this large, because the array is sorted ascending and the left pointer
only ever moves right. So that right-hand value cannot be part of any
solution, and discarding it loses nothing. The symmetric argument handles the
too-small case.

Each step eliminates one element, so the search finishes in at most $n$ steps
--- the same answer a hash map would give, with no extra memory.

\solheading
\begin{lstlisting}[style=java]
public static int[] twoSumSorted(int[] a, int target) {
    if (a == null) return null;
    int i = 0, j = a.length - 1;
    while (i < j) {
        long sum = (long) a[i] + a[j];
        if (sum == target) return new int[] {i, j};
        if (sum < target) i++; else j--;
    }
    return null;
}
\end{lstlisting}

On an \emph{unsorted} array this does not apply --- use a \texttt{HashMap}
from value to index instead, also $O(n)$.

\cost{$O(n)$ time, $O(1)$ space.}
\trap{\texttt{a[i] + a[j]} overflows \texttt{int} near the limits. Widen to
\texttt{long} before adding, as above.}

% ------------------------------------------------------------------
\problem{Maximum subarray sum}{Medium}

Return the largest sum obtainable from a contiguous, non-empty subarray.

\noindent\texttt{[-2,1,-3,4,-1,2,1,-5,4]} gives 6.

\idea{Kadane. At each element the best subarray ending here is either this
element alone, or this element extending the best one ending just before it.}

\begin{dia}
  \cellrowst{a}{0}{cell/-2,cell/1,cell/-3,hot/4,hot/-1,hot/2,hot/1,cell/-5,cell/4}
  \idxrow{-0.62}{8}
  \node[note, anchor=east] at (-0.5,0) {input};
  \cellrowst{e}{-1.5}{cell/-2,cell/1,cell/-2,soft/4,soft/3,soft/5,soft/6,cell/1,cell/5}
  \node[note, anchor=east] at (-0.5,-1.5) {best ending here};
  \draw[flow] (4.8,-1.15) -- (4.8,-0.9);
  \caplabel{4.8}{-2.15}{maximum is 6};
\end{dia}

\textbf{How it works.} The second row is the whole algorithm. At each index
you only ever ask one question: is it better to start fresh here, or to carry
forward what came before? If the running total so far is negative it can only
hurt, so you drop it --- that is exactly what
\texttt{Math.max(a[i], endingHere + a[i])} encodes.

Notice that the best overall answer is not necessarily the last value of the
running total; it is the maximum ever reached. That is why \texttt{best} is
tracked separately.

\solheading
\begin{lstlisting}[style=java]
public static int maxSubarraySum(int[] a) {
    if (a == null || a.length == 0) return 0;
    long best = a[0], endingHere = a[0];
    for (int i = 1; i < a.length; i++) {
        endingHere = Math.max(a[i], endingHere + a[i]);
        best = Math.max(best, endingHere);
    }
    return (int) best;
}
\end{lstlisting}

\cost{$O(n)$ time, $O(1)$ space.}
\trap{Initialising \texttt{best} to \texttt{0} returns 0 on an all-negative
array, when the answer is the least negative element. The statement says
non-empty, so the empty subarray is not an option.}

% ------------------------------------------------------------------
\problem{Rotate an array by k}{Medium}

Shift every element right by $k$ positions, wrapping around, in place.

\idea{Reverse the whole array, then reverse the first $k$ and the rest
separately.}

\begin{dia}
  \node[note, anchor=east] at (-0.5,0.9) {start};
  \cellrowst{a}{0.9}{cell/1,cell/2,cell/3,cell/4,cell/5}
  \node[note, anchor=east] at (-0.5,0) {reverse all};
  \cellrowst{b}{0}{soft/5,soft/4,soft/3,soft/2,soft/1}
  \node[note, anchor=east] at (-0.5,-0.9) {reverse each part};
  \cellrowst{c}{-0.9}{hot/4,hot/5,warm/1,warm/2,warm/3}
  \draw[span] (-0.4,-1.5) -- (1.2,-1.5);
  \draw[span] (1.4,-1.5) -- (3.6,-1.5);
  \caplabel{0.4}{-2.0}{first $k=2$};
  \caplabel{2.5}{-2.0}{the rest};
\end{dia}

\textbf{How it works.} Rotating right by $k$ means the last $k$ elements move
to the front, keeping their order, and the rest follow, keeping theirs.
Reversing the whole array brings the last $k$ to the front but in the wrong
order, and does the same to the rest. Reversing each of those two blocks
individually puts both back in order. Three linear reversals, no second
array.

The normalisation \texttt{k = ((k \% n) + n) \% n} handles two cases at once:
$k$ larger than the array, and negative $k$. A single \texttt{\%} leaves a
negative result negative.

\solheading
\begin{lstlisting}[style=java]
public static void rotate(int[] a, int k) {
    if (a == null || a.length == 0) return;
    k = ((k % a.length) + a.length) % a.length;   // handles negative k
    reverse(a, 0, a.length - 1);
    reverse(a, 0, k - 1);
    reverse(a, k, a.length - 1);
}

private static void reverse(int[] a, int i, int j) {
    while (i < j) {
        int t = a[i]; a[i] = a[j]; a[j] = t;
        i++; j--;
    }
}
\end{lstlisting}

\cost{$O(n)$ time, $O(1)$ space.}
\trap{Not normalising $k$. With \texttt{k = 7} on a length-5 array,
\texttt{reverse(a, 0, 6)} indexes past the end and throws.}

% ------------------------------------------------------------------
\problem{Merge two sorted arrays}{Easy}

Both inputs are sorted ascending. Return one sorted array with all the
elements.

\idea{One index into each; repeatedly take the smaller head, then drain
whichever array still has elements.}

\begin{dia}
  \node[note, anchor=east] at (-0.5,0.9) {a};
  \cellrowst{a}{0.9}{dead/1,soft/5,cell/9}
  \node[note, anchor=east] at (-0.5,0) {b};
  \cellrowst{b}{0}{dead/2,soft/3,cell/8}
  \node[ptr] at (0.8,1.5) {i};
  \node[ptr] at (0.8,-0.75) {j};
  \node[note, anchor=east] at (-0.5,-1.6) {out};
  \cellrowst{o}{-1.6}{hot/1,hot/2,cell/{},cell/{},cell/{},cell/{}}
  \node[note, anchor=west] at (4.2,-1.6) {next: compare 5 and 3};
\end{dia}

\textbf{How it works.} At every step the smallest element not yet taken must
be at the head of one of the two arrays, because each is sorted. So one
comparison is enough to decide what goes next --- there is no searching.

The two drain loops after the main one are not optional. The main loop exits
as soon as \emph{either} array runs out, and whatever remains in the other is
already sorted and larger than everything written so far, so it can be copied
straight across.

Using \texttt{<=} rather than \texttt{<} in the comparison keeps the merge
stable: equal elements from \texttt{a} come out before those from \texttt{b}.

\solheading
\begin{lstlisting}[style=java]
public static int[] merge(int[] a, int[] b) {
    if (a == null) a = new int[0];
    if (b == null) b = new int[0];
    int[] out = new int[a.length + b.length];
    int i = 0, j = 0, k = 0;
    while (i < a.length && j < b.length)
        out[k++] = (a[i] <= b[j]) ? a[i++] : b[j++];
    while (i < a.length) out[k++] = a[i++];
    while (j < b.length) out[k++] = b[j++];
    return out;
}
\end{lstlisting}

\cost{$O(n+m)$ time and space.}
\trap{Forgetting the drain loops, which silently truncates the output.}

% ------------------------------------------------------------------
\problem{Product of the array except self}{Medium}

Position $i$ of the result holds the product of every element except
\texttt{a[i]}. No division allowed.

\idea{Two sweeps. Left to right, store the product of everything \emph{before}
$i$. Right to left, multiply in the product of everything \emph{after} $i$,
carried in a single variable.}

\begin{dia}
  \node[note, anchor=east] at (-0.5,0.9) {input};
  \cellrowst{a}{0.9}{cell/1,cell/2,cell/3,cell/4}
  \node[note, anchor=east] at (-0.5,0) {prefix};
  \cellrowst{p}{0}{soft/1,soft/1,soft/2,soft/6}
  \draw[flow] (-0.1,0.45) -- (2.5,0.45);
  \node[note, anchor=east] at (-0.5,-0.9) {suffix};
  \cellrowst{s}{-0.9}{warm/24,warm/12,warm/4,warm/1}
  \draw[flow] (2.5,-0.45) -- (-0.1,-0.45);
  \node[note, anchor=east] at (-0.5,-1.8) {product};
  \cellrowst{r}{-1.8}{hot/24,hot/12,hot/8,hot/6}
\end{dia}

\textbf{How it works.} The product of everything except position $i$ splits
cleanly into two halves: what is to its left and what is to its right. Each
half is a running product, computable in one sweep.

The implementation is tighter than the diagram suggests. The first loop writes
the prefix products directly into the output array; the second loop walks
backwards, multiplying each slot by a running suffix carried in one variable.
So only the output array is allocated, and the extra space is $O(1)$.

\solheading
\begin{lstlisting}[style=java]
public static int[] productExceptSelf(int[] a) {
    if (a == null || a.length == 0) return new int[0];
    int n = a.length;
    int[] out = new int[n];
    out[0] = 1;
    for (int i = 1; i < n; i++) out[i] = out[i - 1] * a[i - 1];
    int suffix = 1;
    for (int i = n - 1; i >= 0; i--) {
        out[i] *= suffix;
        suffix *= a[i];
    }
    return out;
}
\end{lstlisting}

\cost{$O(n)$ time, $O(1)$ extra space beyond the output.}
\trap{The division shortcut --- total product divided by \texttt{a[i]} ---
breaks on one zero and is undefined on two. That is exactly why the statement
forbids it.}

% ------------------------------------------------------------------
\problem{Container with most water}{Medium}

Each element is the height of a vertical line. Pick two so the rectangle
between them holds the most water: the shorter height times the distance.

\idea{Start at the widest pair and close in, always moving the \emph{shorter}
line.}

\begin{dia}
  \foreach \x/\h in {0/1.8, 0.8/8, 1.6/6, 2.4/2, 3.2/5, 4.0/4, 4.8/8, 5.6/3, 6.4/7} {
    \draw[edge, line width=2pt] (\x,0) -- (\x,\h*0.17);}
  \draw[edge] (-0.3,0) -- (6.7,0);
  \fill[dcool, opacity=0.65] (0.8,0) rectangle (4.8,1.36);
  \draw[flow] (0.8,1.75) -- (0.8,1.45);
  \draw[flow] (4.8,1.75) -- (4.8,1.45);
  \node[ptr] at (0.8,1.95) {i};
  \node[ptr] at (4.8,1.95) {j};
  \caplabel{2.8}{-0.45}{area = $\min(8,8)\times 5 = 40$};
\end{dia}

\textbf{How it works.} The area is limited by the shorter of the two lines, so
think about what happens when you move each one inward. Moving the
\emph{taller} line loses width and cannot raise the limiting height --- the
shorter line is still there, still capping the area. So no better rectangle
uses that pair of positions, and the move is strictly wasteful.

Moving the shorter line, by contrast, at least gives the height a chance to
increase. That asymmetry is what makes it safe to discard one candidate per
step and finish in a single pass.

\solheading
\begin{lstlisting}[style=java]
public static int maxArea(int[] height) {
    if (height == null || height.length < 2) return 0;
    int i = 0, j = height.length - 1, best = 0;
    while (i < j) {
        int h = Math.min(height[i], height[j]);
        best = Math.max(best, h * (j - i));
        if (height[i] < height[j]) i++; else j--;
    }
    return best;
}
\end{lstlisting}

\cost{$O(n)$ time, $O(1)$ space.}
\trap{Moving whichever pointer you feel like, or both at once. The proof that
nothing is missed depends on always discarding the shorter line.}
